\documentclass[twocolumn]{aastex631}
\begin{document}
\title{The reionization ionizing-photon-budget shortfall for star-forming galaxies is not robust to systematics at z\~{}7 (o32 systematic corner)}
\author{NebulaMind Lab (autonomous pipeline)}
\affiliation{NebulaMind Lab, \url{https://lab.nebulamind.net}}
\begin{abstract}
Reionization ionizing-photon-budget reconciliation at z\~{}7: star-forming galaxies require f\_esc=0.105 (+0.106/-0.054) to close the budget (Madau-Dickinson SFRD, log xi\_ion=25.5+/-0.15, clumping C=2-5, JWST-SFRD tail), versus indirect-proxy-inferred f\_esc=0.080 (+0.146/-0.051) from LzLCS O32/beta calibrations. Median delta(required-inferred)=+0.020 dex-frac (16-84\%: -0.119 to +0.130); 58\% of the systematic MC shows a shortfall. Conclusion: the budget CLOSES within the systematic, and the sign holds under both O32 and beta calibrations. The result is bounded by the xi\_ion x clumping x proxy-calibration systematic, not by statistics. Generated autonomously from public data (jwst) via the NebulaMind Lab runner: a bounded, reproducible, descriptive study.
\end{abstract}
\section{Introduction}
Reconciling the reionization ionizing-photon-budget has been a longstanding challenge in understanding the early universe. Recent studies have highlighted potential discrepancies between the observed star formation rate density (SFRD) and the required ionizing photon production to drive reionization [Muñoz2024]. This issue is further complicated by uncertainties in key parameters such as the escape fraction of ionizing photons (f\_esc), the ionizing efficiency of galaxies (xi\_ion), and the clumping factor of the intergalactic medium (C) [Davies2021, Park2022].
\section{Data and method}
To address this challenge, we adopt a literature-anchored budget calculation approach. We use the cosmic SFRD from Madau \& Dickinson's (2014) analytic fitting function, along with published values for xi\_ion and f\_esc proxy calibrations from LzLCS [Chisholm+22, Flury+22] and Simmonds+24. This method allows us to systematically reconcile the reionization photon budget without relying on new observational data.
\section{Result}
Our calculation reveals that star-forming galaxies at z\~{}7 require an escape fraction of f\_esc=0.105 (+0.106/-0.054) to close the ionizing-photon-budget, assuming a Madau-Dickinson SFRD, log xi\_ion=25.5±0.15, and clumping factor C between 2-5. In contrast, indirect-proxy-inferred f\_esc values from LzLCS O32/beta calibrations yield f\_esc=0.080 (+0.146/-0.051). The median difference between the required and inferred escape fractions is +0.020 dex-frac (16-84\%: -0.119 to +0.130), with 58\% of systematic Monte Carlo simulations showing a shortfall.
\begin{figure}\centering\includegraphics[width=\columnwidth]{result.png}\caption{The reionization ionizing-photon-budget shortfall for star-forming galaxies is not robust to systematics at z\~{}7 (o32 systematic corner)}\end{figure}
\section{Caveats}
However, it is essential to acknowledge the limitations of our approach. Our calculation relies on automated, single-selection, uncalibrated measurements, which may introduce biases and uncertainties. The systematic errors in xi\_ion x clumping x proxy-calibration dominate the result, rather than statistical uncertainties. Additionally, the use of published values for f\_esc proxy calibrations assumes that these calibrations are accurate and applicable to our specific context, which may not be entirely valid. Further research is needed to refine these parameters and improve the accuracy of reionization photon budget calculations.
\section*{References}
\footnotesize
[Muoz2024] Muñoz et al. (2024). Reionization after JWST: a photon budget crisis?. bibcode 2024MNRAS.535L..37M \par [Park2022] Park et al. (2022). Calibrating excursion set reionization models to approximately conserve ionizing photons. bibcode 2022MNRAS.517..192P \par [Duncan2015] Duncan et al. (2015). Powering reionization: assessing the galaxy ionizing photon budget at z < 10. bibcode 2015MNRAS.451.2030D \par [Davies2021] Davies et al. (2021). The Predicament of Absorption-dominated Reionization: Increased Demands on Ionizing Sources. bibcode 2021ApJ...918L..35D \par [Madau2017] Madau (2017). Cosmic Reionization after Planck and before JWST: An Analytic Approach. bibcode 2017ApJ...851...50M

\end{document}
